[ C ] The one-sided z-test on the mean, option A, fails because the sparse signal is overwhelmed by noise. The expected number of signals is n epsilon_n = n^{1-beta} = n^{0.15}, which grows, but mu_n = sqrt(2 0.45 log n) = sqrt(0.9 log n). The mean is E approx epsilon_n mu_n = n^{-beta} sqrt(0.9 log n) = n^{-0.85} (log n)^{0.5}, which goes to zero, so sqrt(n) bar X has variance roughly 1 plus negligible signal, yielding power bounded below 0.05 asymptotically. Option B, rejecting for large max X_i, also fails. Under H0, max X_i / sqrt(2 log n) -> 1 in probability by extreme value theory for Gaussians. Under H1, the signal strength mu_n / sqrt(2 log n) = sqrt(0.45) < 1, so even the strongest signals are o_p(sqrt(2 log n)), indistinguishable from null maxima asymptotically. The Higher Criticism test, option C, succeeds. It rejects for sup_t [ (# {i: X_i > t} / n - Phi(-t) ) / sqrt(Phi(-t)(1-Phi(-t))/n ) ] exceeding a threshold near sqrt(2 log log n). At t approx mu_n = sqrt(0.9 log n), the number of exceedances is Bin(n, epsilon_n Phi(- (t - mu_n))) approx Poisson(lambda_n) with lambda_n approx n^{0.15} (1/sqrt(2pi)) >> log log n. The HC statistic at this t exceeds sqrt(2 log log n) -> infinity, yielding power -> 1 by Donoho-Jin theory for this sparse regime (beta < 3/4 would be denser, but beta=0.85 fits their detection boundary for r=0.45 < 0.5).
